
Self-supervised learning (SSL) models trained via contrastive methods such as SimCLR have demonstrated the capacity to achieve high worst-group accuracy (WGA) on spurious correlation benchmarks. Mehta et al. (2022) reported that frozen ViT-H-14 embeddings combined with linear probes reached 90.13\% WGA on the Waterbirds dataset without access to group labels during training. The geometric properties of SSL embeddings that enable this fairness performance are not fully understood.

Spurious correlations occur when a feature correlated with the target label in the training distribution does not have a causal relationship with the true task. In the Waterbirds dataset, background type (land vs. water) exhibits 93\% correlation with bird type (landbird vs. waterbird) in the training split, but the background is not causally relevant to bird classification. Models that rely on this spurious correlation perform poorly on minority groups where the correlation is violated.

One explanation for SSL's fairness performance is that spurious features form discrete, geometrically separable clusters in the embedding space. Under this hypothesis, the InfoNCE contrastive loss used in SimCLR would pull together samples sharing spurious features (e.g., water backgrounds), creating dense clusters. If this cluster structure exists and is measurable via metrics such as Adjusted Mutual Information (AMI), it could serve as a diagnostic tool to identify when cluster-based fairness interventions will be effective. Additionally, methods such as LA-SSL \citep{zhu2023la-ssl} that improve fairness through learning-speed aware resampling could potentially operate by dispersing these spurious clusters.

This hypothesis has not been directly tested. Prior work has demonstrated that linear classifiers achieve high WGA on SSL embeddings, but linear separability does not necessarily imply discrete clusterability. Spurious features could form continuous gradients that support linear decision boundaries without forming discrete density modes detectable by clustering algorithms.

We conducted experiments to test three specific mechanism hypotheses:

\textbf{M1 (Cluster Formation):} Standard SSL training with InfoNCE loss creates geometrically separable spurious clusters with AMI $\geq$ 0.4.

\textbf{M2 (Diagnostic Power):} Clusterability (AMI) predicts the efficacy of cluster-balanced retraining interventions, such that models with AMI $\geq$ 0.4 show WGA improvement $\geq$ 2 percentage points, while models with AMI < 0.3 show < 0.5 percentage points improvement.

\textbf{M3 (Geometric Dispersion):} LA-SSL reduces clusterability by $\geq$ 30\% compared to standard SimCLR while preserving linear separability ($\Delta AUC$ < 0.05).

Our experiments were conducted as proof-of-concept validation with 20-epoch training runs prior to committing resources to full-scale 100-epoch experiments. Under these conditions, all three mechanism hypotheses were not supported. AMI values remained below the 0.4 threshold, LA-SSL increased rather than decreased AMI, and linear separability remained high despite low clusterability.

The primary contributions of this work are: (1) empirical measurements of clusterability in SSL embeddings trained on a spurious correlation dataset, (2) demonstration that linear separability and discrete clusterability are dissociable properties under the tested conditions, and (3) validated implementations of cluster-based retraining mechanisms suitable for future experimental work.
